%% kapa: exam=1 round=31 year=2020 period=1 form=A subject=economics subject_ko=경제학원론 seq=10 number=50
%% source: Q-Net 감정평가사 기출문제 (cst00309) artlSeq=5207219 file=extracted/1008/5207219/A형/A형/2020년도 제31회 감정평가사 1차 시험_1교시_A형.pdf
%% conversion: kordoc 4.24.51 -> tools/gen1.py (auto); review=proofread-sonnet; flags=none
\begin{question}{50}
효용을 극대화하는 甲의 효용함수는 $U(x,y)=min\{x,y\}$이다. 소득이 1,800, $X$재와 $Y$재의 가격은 각각 10 이다. $X$재 가격만 8 로 하락할 때, 옳은 것을 모두 고른 것은? (단, $x$는 $X$재 소비량, $y$는 $Y$재 소비량)
\begin{framedtext}
ㄱ. $X$재 소비량의 변화 중 대체효과는 0이다. \par
ㄴ. $X$재 소비량의 변화 중 소득효과는 10이다. \par
ㄷ. 한계대체율은 하락한다. \par
ㄹ. $X$재 소비는 증가하고 $Y$재 소비는 감소한다. \par
\end{framedtext}
\begin{choices}
  \item ㄱ, ㄴ
  \item ㄱ, ㄷ
  \item ㄴ, ㄷ
  \item ㄴ, ㄹ
  \item ㄷ, ㄹ
\end{choices}
\end{question}
\answer{1}
